\documentclass[11pt]{article}

\usepackage[T1]{fontenc}
\usepackage[utf8]{inputenc}
\usepackage{lmodern}
\usepackage[a4paper,margin=2.6cm]{geometry}
\usepackage{amsmath}
\usepackage{amssymb}
\usepackage{booktabs}
\usepackage{array}
\usepackage{microtype}
\usepackage{listings}
\usepackage{xcolor}
\usepackage[htt]{hyphenat}
\usepackage[colorlinks=true,linkcolor=black,citecolor=black,urlcolor=blue!50!black]{hyperref}
\usepackage{xurl}

\definecolor{codegray}{gray}{0.45}
\definecolor{codebg}{gray}{0.97}
\lstset{
  basicstyle=\ttfamily\small,
  backgroundcolor=\color{codebg},
  commentstyle=\color{codegray},
  keywordstyle=\bfseries,
  breaklines=true,
  columns=fullflexible,
  keepspaces=true,
  frame=single,
  framerule=0.2pt,
  xleftmargin=1em,
  xrightmargin=1em,
}

\title{Stigmergic-Quorum Admission Control: A Practitioner's Guide}
\author{\texttt{moderation-pipeline-prototype}\thanks{Companion to the
    paper \emph{A Decay-Weighted, Source-Independence-Discounted,
    Load-Adaptive Admission Gate for NP-Complete Argumentation Stages}
    (same package), which carries the formal model, the two
    propositions and their proofs, the calibration-tension corollary,
    and the prior-art survey. This guide carries the applied material:
    walkthrough, tuning procedure, troubleshooting, adaptation.}}
\date{August 2026}

\begin{document}
\maketitle

\section{When you need this}
\label{sec:when}

All three conditions together:

\begin{enumerate}
\item Your system has a cheap, tractable fast path and an expensive,
  worst-case-exponential (or otherwise costly) fallback path, and cases
  fall into the fallback whenever the fast path returns something
  uninformative.
\item Admission to the fallback is currently gated by nothing more than
  ``has this case received at least one confirming signal'' --- the
  first plausible-looking report, argument, or vote lets it through.
\item An adversary can manufacture that one confirming signal cheaply,
  in bulk, in a single burst, without needing to sustain anything or
  coordinate identities over time.
\end{enumerate}

The reference deployment is a four-stage content-moderation pipeline
whose Stage~2b falls back to NP-complete preferred-extension enumeration
in an abstract argumentation framework (Dung, 1995) when the grounded
extension is empty and a case's escalation timeout has expired. Nothing
below is specific to argumentation frameworks, though: the gate treats
the expensive stage as an opaque resource pool with a release signal,
and slots into any admission boundary matching the three conditions
above (job queues that fall back to an exhaustive solver, moderation
escalations to human review, rate-limited access to an expensive model
call).

\medskip
\noindent
\textbf{What you get.} A burst of any size, in a single tick, from any
number of distinct or shared identities, is blocked unconditionally.
A sustained, patient, identity-diverse attacker is not blocked, but is
delayed by roughly a factor of five under the reference tuning ---
\S\ref{sec:tuning} shows how that factor is a dial, not a constant, and
what it costs to turn it.

\section{The gate in three rules}
\label{sec:rules}

\begin{enumerate}
\item \textbf{Decay every tick, before anything else.} Each case carries
  a score that shrinks by a fixed fraction $\rho$ every discrete time
  step, whether or not the case received any reinforcement that step.
  This is the ordinary pheromone-evaporation update
  ($\tau \leftarrow (1-\rho)\tau$) from ant colony optimization, applied
  to an admission-pressure score instead of a trail strength. Decay runs
  \emph{first}, so a case that goes quiet starts losing ground
  immediately, not after some grace period.
\item \textbf{Cap the tick's contribution, and discount repeats from the
  same source.} Within one tick, dedupe the reinforcing sources first
  (a source submitting five items in one tick is one reinforcement
  event, not five), then sum: 1.0 per source seen for the first time in
  the last $w$ ticks, a small fraction $\delta$ per source seen again
  inside that window. Cap the sum at a fixed ceiling $\gamma$
  \emph{after} the dedupe-and-discount step, never before. Getting this
  order wrong is the single most consequential implementation mistake in
  this whole gate --- \S\ref{sec:pitfalls} covers it as the first
  troubleshooting entry, because it happened here.
\item \textbf{Raise the bar when the expensive tier is busy.} The
  admission threshold is not a constant; it is
  $\theta_0(1 + \kappa L/C)$, where $L$ is the number of cases currently
  occupying the expensive tier, system-wide, right now --- not a
  cumulative count of everything ever admitted. Getting \emph{that}
  distinction wrong is the second troubleshooting entry.
\end{enumerate}

\section{Annotated walkthrough}
\label{sec:walkthrough}

The reference implementation is \texttt{src/admission\_control.py},
class \texttt{StigmergicQuorumAdmissionControl}. Everything below is
that file, in reading order, with the reasoning moved out of the
comments and into prose.

\subsection{The tunables and their defaults}

\begin{lstlisting}
DEFAULT_DECAY_RATE = 0.10           # rho
DEFAULT_TICK_CAP = 1.0              # gamma
DEFAULT_INDEPENDENCE_WINDOW = 3     # w, in ticks
DEFAULT_REPEAT_DISCOUNT = 0.05      # delta
DEFAULT_BASE_THRESHOLD = 4.0        # theta_0
DEFAULT_LOAD_SENSITIVITY = 0.5      # kappa
DEFAULT_CAPACITY = 5                # C
\end{lstlisting}

These seven numbers are not independent knobs to be tuned one at a time
against a vibe. The paper's Corollary~1 (the calibration tension) proves
that the ratio $x = \theta_0\rho/\gamma$, evaluated at zero load, governs
\emph{both} how easy a genuine, moderately-corroborated dispute is to
admit \emph{and} how many sustained ticks a fully saturated attacker
needs. There is no way to move $x$ that helps one side without hurting
the other. The reference values above were chosen by measuring that
trade-off directly (not tuned backward from a desired test result): they
require an attacker to sustain roughly five ticks of full-saturation
effort, while keeping the bar clearable for a dispute that gets an
independent reinforcement on the order of 40--50\% of ticks on average.
If your deployment's genuine-dispute reinforcement rate is different,
re-derive $x$ for your population before touching any single constant
(\S\ref{sec:tuning}).

\subsection{The per-case state and the tick loop}

\begin{lstlisting}
@dataclass
class CaseAdmissionState:
    score: float = 0.0
    last_reinforced_tick: dict[str, int] = field(default_factory=dict)
    admitted: bool = False
    admitted_tick: int | None = None
\end{lstlisting}

\texttt{last\_reinforced\_tick} is keyed by cohort id, not by event: it
is the mechanism's entire memory of who has recently reinforced this
case, and it is what makes the independence discount possible without
storing full event history.

\begin{lstlisting}
def tick(self, t: int, reinforcements: dict[str, list[str]]) -> set[str]:
    newly_admitted: set[str] = set()
    touched = set(self._cases) | set(reinforcements)
    for case_id in touched:
        state = self._cases.setdefault(case_id, CaseAdmissionState())
        if state.admitted:
            continue
        # (1) decay first
        state.score *= (1 - self.decay_rate)
        # (2) accumulate this tick's reinforcement, deduped, weighted,
        #     capped
        contrib = 0.0
        for cohort in set(reinforcements.get(case_id, [])):
            last = state.last_reinforced_tick.get(cohort)
            fresh = last is None or (t - last) >= self.independence_window
            contrib += 1.0 if fresh else self.repeat_discount
            state.last_reinforced_tick[cohort] = t
        contrib = min(contrib, self.tick_cap)
        state.score += contrib
        # (3) admit only if the load-adaptive threshold is crossed
        if state.score >= self.current_threshold():
            state.admitted = True
            state.admitted_tick = t
            self._pending_admitted.add(case_id)
            newly_admitted.add(case_id)
    return newly_admitted
\end{lstlisting}

Notice \texttt{set(reinforcements.get(case\_id, []))} on the cohort loop
--- that \texttt{set()} call is the dedupe step, and it runs before the
discount logic, not after. This is the fix for the first bug in
\S\ref{sec:pitfalls}; if you are adapting this gate to a new language or
rewriting the loop, put the dedupe there and nowhere later.

\begin{lstlisting}
def current_threshold(self) -> float:
    load = len(self._pending_admitted)
    return self.base_threshold * (1 + self.load_sensitivity * load / self.capacity)

def release(self, case_id: str) -> None:
    self._pending_admitted.discard(case_id)
\end{lstlisting}

\texttt{release()} is not optional plumbing you can skip in a first
integration. \S\ref{sec:pitfalls}'s second entry is what happens when you
do.

\subsection{The closed-form helper}

\begin{lstlisting}
def min_ticks_to_admission(threshold, tick_cap, decay_rate) -> float:
    ceiling_fraction = threshold * decay_rate / tick_cap
    if ceiling_fraction >= 1.0:
        return math.inf
    return math.log(1 - ceiling_fraction) / math.log(1 - decay_rate)
\end{lstlisting}

This function is not a side utility; it is Proposition~2 of the
companion paper made executable, and it is used twice: once to
\emph{design} \texttt{base\_threshold} against a target attacker cost,
and once inside the adversarial test suite to check that the simulated
gate actually admits a saturated attacker at the tick the closed form
predicts. If you change any of \texttt{decay\_rate}, \texttt{tick\_cap},
or \texttt{base\_threshold}, re-run this function before re-running the
simulation --- if the two disagree by more than a tick, something in the
implementation has drifted from the model it is supposed to realize.

\section{Tuning procedure}
\label{sec:tuning}

Do not hand-tune the seven constants independently. Follow this order:

\begin{enumerate}
\item \textbf{Fix $\rho$ from your tick granularity.} $\rho$ sets the
  half-life of unreinforced pressure, $\ln 2/\rho$ ticks. Pick a tick
  length matching how often your system can plausibly observe
  reinforcement events (minutes, not milliseconds, for a
  human-moderation pipeline), then pick $\rho$ so the half-life matches
  how long a genuinely urgent dispute should stay ``warm'' between
  independent reinforcements.
\item \textbf{Fix $\gamma$ and $\delta$ from your identity-diversity
  ceiling.} $\gamma=1$ (a single fresh source per tick fully saturates
  the cap) is the conservative default. $\delta$ should be small enough
  that your expected same-cohort repeat rate cannot reach the threshold
  even given unlimited time --- check this with
  Proposition~1: the ceiling is $\delta\gamma/\rho$, and it must sit
  below $\theta_0$ by a margin you are comfortable with, not just barely
  under it.
\item \textbf{Fix $\theta_0$ from the calibration tension, explicitly.}
  Compute $x = \theta_0\rho/\gamma$ for candidate values of $\theta_0$
  and, for each, compute both (a) the genuine-case admission probability
  a moderately corroborated real dispute in your population is likely to
  clear, and (b) $n_{\min}(x)$ from \texttt{min\_ticks\_to\_admission}.
  Pick the value where (b) is large enough to matter operationally
  (multiple sustained rounds of attacker effort, not one) while (a) is
  still comfortably clearable by your real traffic. There is no value
  that maximizes both; Corollary~1 proves that formally, so do not spend
  time looking for one.
\item \textbf{Fix $\kappa$ and $C$ last, against your actual expensive-tier
  capacity.} $C$ should be your real concurrent capacity for the
  expensive stage (however many cases it can process at once without
  degrading), and $\kappa$ controls how sharply the bar rises as that
  capacity fills. Start at $\kappa=0.5$ and increase it if the ablation
  test (\S\ref{sec:tuning}, and the paper's Attack~5) does not spread
  concurrent admissions out enough for your operational tolerance.
\item \textbf{Re-run the adversarial test suite with your new constants,
  not just the sanity control.} A tuning that only checks that genuine
  cases still get admitted, without re-running the burst and
  sustained-flood attacks, has not actually been validated.
\end{enumerate}

\section{Integration checklist}
\label{sec:checklist}

\begin{itemize}
\item[$\square$] Confirmed the admission boundary you are protecting
  matches the three conditions of \S\ref{sec:when} --- in particular,
  that ``admission'' is currently a single-signal, single-tick decision,
  not something already rate-limited some other way.
\item[$\square$] Identified your real tick unit and your real
  reinforcement-event source (what, concretely, counts as ``a cohort
  reinforcing a case'' in your system --- a report, a rebuttal, a
  resubmission).
\item[$\square$] Chosen a cohort-identity proxy and written down its
  honest limitation. This prototype uses a structural
  \texttt{rater\_creation\_cohort} id; whatever proxy you use, the gate's
  independence discount is only as good as that proxy's resistance to
  spoofing (\S\ref{sec:limits}).
\item[$\square$] Wired \texttt{release()} into whatever event marks a
  case as finished in the expensive tier --- not into the admission
  event, and not left uncalled (\S\ref{sec:pitfalls}, second entry).
\item[$\square$] Re-derived $x=\theta_0\rho/\gamma$ for your own
  genuine-case reinforcement rate rather than copying the reference
  constants verbatim (\S\ref{sec:tuning}).
\item[$\square$] Re-run all six tests in
  \texttt{tests/test\_admission\_control.py} against your constants, and
  read the printed numbers, not just the pass/fail status --- the
  sanity-control admission rate and the measured attacker delay are the
  numbers that tell you whether the tuning is actually workable for your
  traffic.
\item[$\square$] Decided, explicitly, what happens to a case that never
  crosses the threshold. This gate has no forced-timeout escape hatch of
  its own; if your system needs one (a human ombudsman review after $N$
  ticks of stalled pressure, say), it is a separate mechanism layered on
  top, not a hidden default here.
\end{itemize}

\section{Troubleshooting}
\label{sec:pitfalls}

Both rows below are real bugs caught while building and adversarially
testing this gate, not hypothetical failure modes.

\begin{center}\small
\begin{tabular}{>{\raggedright\arraybackslash}p{4.3cm}>{\raggedright\arraybackslash}p{4.3cm}>{\raggedright\arraybackslash}p{4.9cm}}
\toprule
Symptom & Root cause & Fix \\
\midrule
A same-cohort flood still reaches the tick cap and gets admitted, even
though the repeat discount is configured correctly &
Same-cohort contributions were summed \emph{before} deduplication within
a tick, so enough discounted entries from one source could still add up
to the cap through sheer in-tick volume &
Dedupe cohorts per tick with a plain \texttt{set()} call
\emph{before} applying the discount and the cap, not after \\
\addlinespace
Under a sustained multi-campaign attack simulation, the admission
threshold climbs steadily and later, genuinely well-corroborated cases
stop getting admitted at all &
\texttt{release()} was never called, so \texttt{\_pending\_admitted}
only ever grows, and \texttt{current\_threshold()} rises monotonically
as a pure artifact of the simulation, not a real property of the
mechanism &
Model a bounded expensive-tier dwell time and call \texttt{release()}
when it elapses; do not loosen the test's assertion to paper over a
threshold that should have come back down \\
\bottomrule
\end{tabular}
\end{center}

The general lesson behind both rows: everything this gate does is
sensitive to \emph{when}, exactly, a mutating step runs relative to the
cap or the load signal it feeds. Whenever you touch the tick loop,
re-check the ordering, not just the arithmetic.

\section{Adapting to a different admission boundary}
\label{sec:adapt}

\begin{enumerate}
\item \textbf{Identify your cohort proxy.} Anything that plausibly
  distinguishes independent sources works: account age plus creation
  cluster, network-origin fingerprint, or (weakest, but sometimes all
  you have) plain user id. The gate's independence discount is only as
  strong as this proxy's resistance to being spoofed fresh every tick
  --- state that limitation explicitly wherever you deploy this, as
  \S\ref{sec:limits} does here.
\item \textbf{Identify your load signal.} It must be a real, currently-occupied
  count that drains as cases leave the expensive tier, not a
  cumulative counter. If your expensive tier has no natural
  ``currently occupying'' notion (e.g.\ a one-shot expensive computation
  rather than a tier with dwell time), the load-adaptive threshold layer
  (\S\ref{sec:rules}, rule 3) may not transfer; the decay-plus-cap layer
  (rules 1--2) still does, on its own.
\item \textbf{Re-derive the calibration tension for your own genuine-case
  distribution.} Do not reuse $\theta_0=4.0$ from a different domain's
  reinforcement pattern; re-run the procedure of \S\ref{sec:tuning}
  against your own traffic.
\item \textbf{Re-run the adversarial suite, adapted to your domain's
  attack shapes.} The five scenarios in the companion paper (single-tick
  burst, persistent same-cohort flood, persistent cohort-spoofed flood,
  spray across many cases, synchronized parallel campaigns) are a
  reasonable starting checklist, not a closed list --- think about
  what a resource-exhaustion attacker in your specific domain would
  actually try first.
\end{enumerate}

\section{Known limits}
\label{sec:limits}

The gate does not close the sustained, identity-diverse attack; it
delays it, by a factor that is itself a tunable trade-off against the
bar for genuine cases (Corollary~1 of the companion paper). It assumes
the cohort-identity proxy has some real cost to spoof; if an attacker
can generate a structurally fresh-looking identity at zero marginal cost
every tick, the delay factor is the entire defense, not a backstop on
top of something stronger. The load-adaptive layer assumes the load
signal is wired to real occupancy, per the second troubleshooting
row above --- get that wrong and the layer's behavior is undefined by
this document. And the gate is an admission control, not a
content-judgment mechanism: it says nothing about whether a case, once
admitted, is decided correctly by the expensive stage it gates access
to.

For the formal model, the two propositions and their proofs, the
calibration-tension corollary, and the prior-art survey, see the
companion paper in this package.

\end{document}
